%% Shiva Mishra — Professional Resume
%% Compile with: pdflatex main.tex
\documentclass[letterpaper,11pt]{article}

\usepackage[T1]{fontenc}
\usepackage[utf8]{inputenc}
\usepackage{lmodern}
\usepackage[margin=0.75in]{geometry}
\usepackage{hyperref}
\usepackage{titlesec}
\usepackage{enumitem}
\usepackage{xcolor}
\usepackage{microtype}
\usepackage{parskip}

% Colours
\definecolor{accent}{HTML}{1A56DB}

% Section formatting
\titleformat{\section}{\large\bfseries\color{accent}}{}{0em}{}[\titlerule]
\titlespacing{\section}{0pt}{10pt}{4pt}

% No page numbers
\pagestyle{empty}

% Tight lists
\setlist[itemize]{noitemsep, topsep=2pt, leftmargin=1.5em}

% Hyperlink style
\hypersetup{colorlinks=true, urlcolor=accent, linkcolor=accent}

% ------------------------------------------------------------------
% Helper commands
% ------------------------------------------------------------------
\newcommand{\resumeitem}[1]{\item #1}

\newcommand{\entry}[4]{%
  \noindent
  \textbf{#1} \hfill \textit{#2} \\
  \textit{#3} \hfill #4 \par
}

% ------------------------------------------------------------------
\begin{document}

% ── Header ─────────────────────────────────────────────────────────
\begin{center}
  {\Huge\bfseries Shiva Mishra} \\[4pt]
  Lead AI/ML Engineer \\[4pt]
  Hyderabad, India \quad|\quad
  \href{mailto:21shiva12mishra@gmail.com}{21shiva12mishra@gmail.com} \quad|\quad
  \href{https://linkedin.com/in/shivamishra1}{linkedin.com/in/shivamishra1} \quad|\quad
  \href{https://github.com/bydefaultuser}{github.com/bydefaultuser}
\end{center}

% ── Summary ────────────────────────────────────────────────────────
\section{Summary}
Lead AI/ML Engineer with production experience designing and deploying multi-agent LLM systems,
RAG pipelines, and network automation tooling at scale.
Specialises in LangChain/LangGraph orchestration, MCP server development, and AWS cloud infrastructure.
Passionate about building AI systems that are observable, testable, and production-ready.

% ── Experience ─────────────────────────────────────────────────────
\section{Experience}

\entry{Lead AI/ML Engineer}{Keyspark Technologies, Hyderabad}{Full-Time}{2023 -- Present}
\begin{itemize}
  \resumeitem{Architected a multi-agent LLM platform using LangGraph with a supervisor–worker pattern,
              reducing manual triage time for network incidents by 60\%.}
  \resumeitem{Built production MCP servers (FastAPI + SSE transport) exposing internal tools to Claude Desktop
              and custom AI agents across the engineering organisation.}
  \resumeitem{Designed and deployed RAG pipelines using ChromaDB and Supabase for domain-specific
              knowledge retrieval, achieving sub-500 ms p95 query latency.}
  \resumeitem{Integrated Nautobot, Batfish, and NIST NVD APIs into automated network compliance workflows,
              cutting CVE review cycles from days to hours.}
  \resumeitem{Provisioned and maintained AWS infrastructure (ECS, ECR, RDS, S3, CloudWatch) with Terraform
              and GitHub Actions CI/CD; drove 99.9\% uptime across all AI services.}
  \resumeitem{Set up observability with VictoriaMetrics and Telegraf for LLM token usage, latency,
              and error rate dashboards.}
\end{itemize}

\vspace{6pt}
\entry{ML Engineer}{Previous Role}{Full-Time}{2021 -- 2023}
\begin{itemize}
  \resumeitem{Fine-tuned BERT and GPT-2 models for domain-specific NLP classification tasks
              (precision/recall > 0.92 on held-out test sets).}
  \resumeitem{Built data pipelines ingesting 10M+ records/day using PostgreSQL, Redis, and Python
              async workers.}
  \resumeitem{Deployed models on AWS EC2 with Docker; introduced model-versioning workflow that
              reduced rollback time from hours to under 5 minutes.}
\end{itemize}

% ── Skills ─────────────────────────────────────────────────────────
\section{Technical Skills}

\noindent
\begin{tabular}{@{}ll}
  \textbf{AI / ML}         & LLMs, LangChain, LangGraph, MCP, RAG, ChromaDB, PyTorch, TensorFlow, \\
                           & Hugging Face, NLP, Fine-tuning \\[4pt]
  \textbf{Cloud / DevOps}  & AWS (EC2, ECS, ECR, RDS, S3, CloudWatch), Terraform, Docker, \\
                           & GitHub Actions, CI/CD, Linux \\[4pt]
  \textbf{Network Automation} & Nautobot, Batfish, VictoriaMetrics, Telegraf, SNMP, Nginx, Cloudflare \\[4pt]
  \textbf{Databases}       & PostgreSQL, Redis, ChromaDB, Supabase \\[4pt]
  \textbf{Languages}       & Python (Expert), SQL, Bash \\
\end{tabular}

% ── Projects ───────────────────────────────────────────────────────
\section{Open-Source Projects}

\noindent\textbf{\href{https://github.com/bydefaultuser/mcp-server-template}{mcp-server-template}}
\hfill \textit{Python, FastAPI, MCP}\\
Production-ready template for building Model Context Protocol servers with HTTP and SSE transports.

\vspace{4pt}
\noindent\textbf{\href{https://github.com/bydefaultuser/multi-agent-ai-starter}{multi-agent-ai-starter}}
\hfill \textit{Python, LangGraph, LangChain}\\
Boilerplate for supervisor–worker multi-agent systems with researcher, coder, and reviewer agents.

\vspace{4pt}
\noindent\textbf{\href{https://github.com/bydefaultuser/network-automation-utils}{network-automation-utils}}
\hfill \textit{Python, Nautobot, NVD API}\\
Utilities for Nautobot REST integration, Cisco IOS config parsing, and NIST NVD CVE lookup.

% ── Education ──────────────────────────────────────────────────────
\section{Education}

\entry{B.Tech, Computer Science \& Engineering}{University Name, Hyderabad}{}{2017 -- 2021}

\end{document}
